\section{Noether's Theorem for Fields}
\label{sec:noether_theorem_fields}

Noether's theorem converts a continuous variational symmetry into a local conservation law. The theorem is not a separate guess for each theory. It follows from the universal variational identity established in the previous section.

\subsection{Statement of the theorem}

Let

\[
S[\phi]
=
\int_{\Omega}
\mathcal{L}
\bigl(
\phi_a,
\partial_{\mu}\phi_a,
x
\bigr)
\,\mathrm{d}^{4}x
\]

be an action for fields \(\phi_a\). Consider the infinitesimal transformation

\[
x'^{\mu}
=
x^{\mu}
+
\varepsilon\xi^{\mu},
\]

\[
\phi'_a(x')
=
\phi_a(x)
+
\varepsilon\Psi_a.
\]

Define

\[
Q_a
=
\Psi_a
-
\xi^{\nu}\partial_{\nu}\phi_a,
\]

and suppose the action changes by a divergence:

\[
\delta S
=
\varepsilon
\int_{\Omega}
\partial_{\mu}K^{\mu}
\,\mathrm{d}^{4}x.
\]

Then the Noether current is

\[
\boxed{
j^{\mu}
=
\frac{\partial\mathcal{L}}
{\partial(\partial_{\mu}\phi_a)}
Q_a
+
\mathcal{L}\xi^{\mu}
-
K^{\mu}
}.
\]

Its divergence satisfies the off-shell identity

\[
\boxed{
\partial_{\mu}j^{\mu}
=
-\mathcal{E}_aQ_a
},
\]

where

\[
\mathcal{E}_a
=
\frac{\partial\mathcal{L}}
{\partial\phi_a}
-
\partial_{\mu}
\left[
\frac{\partial\mathcal{L}}
{\partial(\partial_{\mu}\phi_a)}
\right].
\]

Therefore, on solutions of the field equations,

\[
\boxed{
\partial_{\mu}j^{\mu}=0
}.
\]

\subsection{Derivation}

The general variational identity is

\[
\delta S
=
\varepsilon
\int_{\Omega}
\left[
\mathcal{E}_aQ_a
+
\partial_{\mu}
\left(
\pi_a^{\mu}Q_a
+
\mathcal{L}\xi^{\mu}
\right)
\right]
\mathrm{d}^{4}x.
\]

For a variational symmetry,

\[
\delta S
=
\varepsilon
\int_{\Omega}
\partial_{\mu}K^{\mu}
\,\mathrm{d}^{4}x.
\]

Subtracting the two expressions gives

\[
\int_{\Omega}
\left[
\mathcal{E}_aQ_a
+
\partial_{\mu}
\left(
\pi_a^{\mu}Q_a
+
\mathcal{L}\xi^{\mu}
-
K^{\mu}
\right)
\right]
\mathrm{d}^{4}x
=
0.
\]

Because the relation holds locally for arbitrary subregions compatible with the transformation,

\[
\mathcal{E}_aQ_a
+
\partial_{\mu}j^{\mu}
=
0.
\]

On shell,

\[
\mathcal{E}_a=0,
\]

so

\[
\partial_{\mu}j^{\mu}=0.
\]

\subsection{Internal-symmetry form}

For a transformation that does not move coordinates,

\[
\xi^{\mu}=0.
\]

If the Lagrangian density is strictly invariant,

\[
K^{\mu}=0.
\]

The current becomes

\[
\boxed{
j^{\mu}
=
\frac{\partial\mathcal{L}}
{\partial(\partial_{\mu}\phi_a)}
\Psi_a
}.
\]

This is the form used for the global phase symmetry of a complex scalar field.

\subsection{Spacetime-symmetry form}

For a scalar field under a coordinate transformation, there may be no intrinsic field change:

\[
\Psi=0.
\]

Then

\[
Q
=
-\xi^{\nu}\partial_{\nu}\phi.
\]

The current is

\[
j^{\mu}
=
-\pi^{\mu}
\xi^{\nu}\partial_{\nu}\phi
+
\mathcal{L}\xi^{\mu}
-
K^{\mu}.
\]

For constant translations, this expression contains the canonical energy-momentum tensor.

\subsection{Quasi-symmetries}

A transformation may change the Lagrangian by a total divergence:

\[
\delta\mathcal{L}
=
\varepsilon\partial_{\mu}K^{\mu}.
\]

Such a transformation is sometimes called a quasi-symmetry. It produces a conserved current just as a strictly invariant Lagrangian does.

The term \(K^{\mu}\) must be included in the current. Omitting it can produce a current whose divergence does not vanish.

\subsection{Worked example: massless scalar shift symmetry}

Take

\[
\mathcal{L}
=
\frac{1}{2}
\partial_{\mu}\phi
\partial^{\mu}\phi.
\]

The transformation is

\[
\phi
\longrightarrow
\phi+\varepsilon.
\]

Thus

\[
\Psi=1,
\qquad
\xi^{\mu}=0,
\qquad
K^{\mu}=0.
\]

The canonical derivative momentum is

\[
\pi^{\mu}
=
\partial^{\mu}\phi.
\]

The Noether current is

\[
\boxed{
j^{\mu}
=
\partial^{\mu}\phi
}.
\]

Its divergence is

\[
\partial_{\mu}j^{\mu}
=
\Box\phi.
\]

The Euler--Lagrange equation is

\[
\Box\phi=0,
\]

so

\[
\boxed{
\partial_{\mu}j^{\mu}=0
}.
\]

\subsection{Explicit breaking of the shift current}

Add a potential:

\[
\mathcal{L}
=
\frac{1}{2}
\partial_{\mu}\phi
\partial^{\mu}\phi
-
V(\phi).
\]

Under the same shift,

\[
\delta\mathcal{L}
=
-\varepsilon V'(\phi).
\]

The transformation is not a symmetry unless \(V'(\phi)\) vanishes identically. The would-be current still satisfies

\[
\partial_{\mu}j^{\mu}
=
\Box\phi.
\]

Using the field equation

\[
\Box\phi+V'(\phi)=0,
\]

we find

\[
\boxed{
\partial_{\mu}j^{\mu}
=
-V'(\phi)
}.
\]

The nonzero divergence measures the explicit breaking of the shift symmetry.

\subsection{Current improvement}

A conserved current is not unique. Let

\[
B^{\mu\nu}
=
-B^{\nu\mu}.
\]

Define

\[
j'^{\mu}
=
j^{\mu}
+
\partial_{\nu}B^{\mu\nu}.
\]

Then

\[
\partial_{\mu}j'^{\mu}
=
\partial_{\mu}j^{\mu}
+
\partial_{\mu}\partial_{\nu}B^{\mu\nu}.
\]

Because partial derivatives commute and \(B^{\mu\nu}\) is antisymmetric,

\[
\partial_{\mu}\partial_{\nu}B^{\mu\nu}=0.
\]

Therefore

\[
\partial_{\mu}j'^{\mu}
=
\partial_{\mu}j^{\mu}.
\]

The added term is called an improvement term. It can change the local appearance of a current without changing its conservation law or, under suitable boundary conditions, its integrated charge.

\subsection{On-shell and off-shell statements}

The identity

\[
\partial_{\mu}j^{\mu}
=
-\mathcal{E}_aQ_a
\]

is off shell: it holds for arbitrary field configurations.

The conservation law

\[
\partial_{\mu}j^{\mu}=0
\]

is on shell: it holds only when the Euler--Lagrange equations are satisfied.

This distinction is important in formal manipulations and in numerical checks. A trial field that does not satisfy the equations of motion need not produce a conserved Noether current.

\subsection{A practical algorithm}

For a proposed continuous symmetry:

\begin{enumerate}
    \item write the infinitesimal coordinate change \(\xi^{\mu}\);
    \item write the intrinsic field change \(\Psi_a\);
    \item form the characteristic \(Q_a=\Psi_a-\xi^{\nu}\partial_{\nu}\phi_a\);
    \item compute the derivative momenta \(\pi_a^{\mu}\);
    \item determine whether the Lagrangian changes by \(\partial_{\mu}K^{\mu}\);
    \item construct
    \[
    j^{\mu}
    =
    \pi_a^{\mu}Q_a
    +
    \mathcal{L}\xi^{\mu}
    -
    K^{\mu};
    \]
    \item compute \(\partial_{\mu}j^{\mu}\);
    \item use the field equations only at the final step.
\end{enumerate}

\subsection{Common mistakes}

Typical errors include:

\begin{enumerate}
    \item identifying a candidate transformation without checking the action;
    \item forgetting the coordinate contribution \(\mathcal{L}\xi^{\mu}\);
    \item omitting the boundary vector \(K^{\mu}\);
    \item claiming conservation off shell;
    \item treating two currents that differ by an improvement as physically distinct without checking their charges;
    \item replacing a local continuity equation by a global statement without specifying boundary conditions.
\end{enumerate}

\subsection{What has been established}

A continuous variational symmetry produces a current

\[
j^{\mu}
=
\pi_a^{\mu}Q_a
+
\mathcal{L}\xi^{\mu}
-
K^{\mu},
\]

whose divergence is proportional to the Euler--Lagrange expressions. On physical solutions,

\[
\partial_{\mu}j^{\mu}=0.
\]

The next section explains how this local equation becomes a conserved charge and what boundary conditions are required.
